\subsection{Normal partial operators and the spectral theorem}

In this section, E is a complex Hilbert space.

DEFINITION 2. --- Let u be a partial operator on E. One says that u is normal if u is closed with dense domain and if its bounded transform b(u) is a normal endomorphism of E.

If \(u \in \mathcal{L}(\mathrm{E})\) this definition coincides with EVT, V, p. 42, def. 4, by virtue of the formulas \(1_{\mathrm{E}} - b(u)^{*}b(u) = \mathrm{W}(u) = (1_{\mathrm{E}} + u^{*}u)^{-1}\) (prop. 1 of IV, p. 262) and \(b(u^{*}) = b(u)^{*}\) (cor. of loc. cit.).

If \(u\) is a self-adjoint partial operator on \(\mathrm{E}\) , then \(b(u)\) is Hermitian (cor. of prop. 2 of IV, p. 264) hence \(u\) is normal.

Let D be the open unit disk in C. We denote by \(\beta\) the function from C to D defined by \(\beta(z)=z/\sqrt{1+|z|^{2}}\) . It is a homeomorphism, whose inverse satisfies \(\beta^{-1}(z)=z/\sqrt{1-|z|^{2}}\) for \(z\in D\) .

Let \(u \in \mathcal{L}(\mathrm{E})\) . It follows from formulas (2) and (3) of proposition 1 of IV, p. 262 that \(u = \beta^{-1}(b(u))\) , and hence

\[
b (u) = \beta (u).\tag{5}
\]

Lemma 2. --- Let X be a locally compact topological space and let μ be a positive measure on X.

\begin{enumerate}
\def\labelenumi{\alph{enumi})}
\item
  Let \(g\) be a function \(\mu\)-measurable on X. The multiplication operator \(m_g\) on \(\mathrm{L}^2 (\mathrm{X},\mu)\) is normal and \(b(m_g) = m_{\beta \circ g}\) ;
\item
  Let \(h\colon X\to D\) be a function \(\mu\)-measurable. The endomorphism \(m_h\) belongs to \(\Omega(L^2(X,\mu))\) and \(B(m_h)=m_{\beta^{-1}\circ h}\) .
\end{enumerate}

The partial operator \(m_g\) is closed with dense domain. Since \(m_g^* m_g = m_{\overline{g}} m_g = m_{|g|^2}\) (prop. 23 of IV, p. 253 and prop. 24 of IV, p. 255) we have

\[
\mathrm{W} (m _ {g}) = - \mathrm{R} (m _ {| g | ^ {2}}, - 1) = m _ {(1 + | g | ^ {2}) ^ {- 1}}
\]

(prop. 22 of IV, p. 252). This implies

\[
b (m _ {g}) = m _ {g} \circ \mathrm{W} (m _ {g}) ^ {1 / 2} = m _ {g} \circ m _ {(1 + | g | ^ {2}) ^ {- 1 / 2}} = m _ {\beta \circ g}
\]

according to formula (2) of IV, p. 262, the cor. of prop. 6 of IV, p. 187 and lemma 6 of IV, p. 254. Finally applying prop. 6 of IV, p. 187, we obtain assertion a).

Let us prove \(b\) ). Since \(1 - |h(x)|^2 > 0\) for all \(x \in \mathrm{X}\) , we have \(m_h \in \Omega(\mathrm{L}^2(\mathrm{X}, \mu))\) (cor. of prop. 6 of IV, p. 187). Since the map \(u \mapsto b(u)\) is injective (prop. 2 of IV, p. 264) and the endomorphisms \(b(\mathrm{B}(m_h))\) and \(b(m_{\beta^{-1}\circ h})\) are equal to \(m_h\) according to assertion \(a\) ), we indeed have \(\mathrm{B}(m_h) = m_{\beta^{-1}\circ h}\) .

THEOREM 1 ("Spectral theorem"). --- Let u be a normal partial operator on the Hilbert space E. There exist a locally compact topological space X, a positive measure μ on X, an isometric isomorphism θ of L²(X,μ) onto E and a continuous function g on X, such that \(u = \theta \circ m_g \circ \theta^{-1}\).

The endomorphism \(b(u)\) of E is normal by definition. By the spectral theorem for normal endomorphisms of E (corollary 1 of IV, p. 193), there therefore exist a locally compact topological space \(\widetilde{X}\), a positive measure \(\mu\) on \(\widetilde{X}\), an isometric isomorphism \(\widetilde{\theta}\) of \(\mathrm{L}^2 (\widetilde{\mathrm{X}},\mu)\) onto E, and a bounded continuous function \(h\) on \(\widetilde{X}\), so that \(b(u)\) coincides with \(\widetilde{\theta}\circ m_h\circ \widetilde{\theta}^{-1}\) .

Let N be the closed subset of \(x \in \widetilde{X}\) such that \(|h(x)| \geqslant 1\) . Since the endomorphism \(1 - b(u)^{*}b(u) = \widetilde{\theta} \circ m_{1-|h|^2} \circ \widetilde{\theta}^{-1}\) is positive and injective, the set N is locally \(\mu\) -negligible (lemma 7 of IV, p. 186 and cor. of prop. 6 of IV, p. 187). We then set X = \(\widetilde{X} - N\) . This is a locally compact space, and the restriction of functions to X induces an isometric isomorphism of \(\mathrm{L}^{2}(\widetilde{\mathrm{X}}, \mu)\) onto \(\mathrm{L}^{2}(\mathrm{X}, \mu|\mathrm{X})\) (prop. 1 of IV, p. 179). The composition of \(\widetilde{\theta}\) and of the restriction of functions to X therefore induces an isometric isomorphism \(\theta\) of \(\mathrm{L}^{2}(\mathrm{X}, \mu|\mathrm{X})\) onto E. Since \(b(u) = \theta \circ m_{h|\mathrm{X}} \circ \theta^{-1}\) , it follows that \(u = \theta \circ m_{\beta^{-1} \circ (h|\mathrm{X})} \circ \theta^{-1}\) (lemma 2). The theorem follows from this formula.

Lemma 3. --- Let u be a normal partial operator on E. We have

\[
\operatorname{Sp} (b (u)) \cap \mathrm{D} = \beta (\operatorname{Sp} (u)).
\]

One has \(b(u) = u \circ \mathrm{W}(u)^{1/2}\) (prop. 1 of IV, p. 262). The endomorphism \(\mathrm{W}(u)^{1/2}\) is injective and its image is the domain of u (formula (1) of IV, p. 262).

Let \(\lambda \in \mathbf{C}\) . The number \(\lambda\) belongs to the resolvent set of \(u\) if and only if \((u - \lambda 1_{\mathrm{E}}) \circ \mathrm{W}(u)^{1/2}\) is a bijective linear map from E onto E (remark 3 of IV, p. 245). Now, by formula (3) of IV, p. 262, we have

\[
(u - \lambda 1 _ {\mathrm{E}}) \circ \mathrm{W} (u) ^ {1 / 2} = b (u) - \lambda (1 _ {\mathrm{E}} - b (u) ^ {*} b (u)) ^ {1 / 2} = f _ {\lambda} (b (u))
\]

where \(f_{\lambda}\) is the continuous function defined on \(\overline{\mathrm{D}}\) by \(z \mapsto z - \lambda (1 - |z|^2)^{1/2}\) . The endomorphism \(f_{\lambda}(b(u))\) is bijective if and only if its spectrum does not contain 0. Since \(\mathrm{Sp}(f_{\lambda}(b(u))) = f_{\lambda}(\mathrm{Sp}(b(u)))\) (cor. 2 of prop. 7 of I, p. 111), this is the case if and only if 0 does not belong to the set \(f_{\lambda}(\mathrm{Sp}(b(u)))\) . Thus, \(\lambda \in \mathrm{Sp}(u)\) if and only if there exists \(z \in \mathrm{Sp}(b(u)) \cap \overline{\mathrm{D}}\) such that

\[
z - \lambda (1 - | z | ^ {2}) ^ {1 / 2} = 0.
\]

This equality, if valid, implies that \(z \in D\) , and means that \(\lambda = \beta^{-1}(z)\) . We conclude that \(\operatorname{Sp}(u) = \beta^{-1}(\operatorname{Sp}(b(u)) \cap D)\) , as stated.

Let u be a partial normal operator on E. Let \(f \in \mathcal{K}(\mathrm{Sp}(u))\) . The map \(z \mapsto f(\beta^{-1}(z))\) from \(\mathrm{Sp}(b(u)) \cap D\) into C is continuous and compactly supported; the unique mapping \(f_{\beta}\) from \(\mathrm{Sp}(b(u))\) into C which extends it by zero is therefore continuous.

DEFINITION 3. --- For every function \(f \in \mathcal{K}(\mathrm{Sp}(u))\) , we define the endomorphism \(f(u)\) of E by \(f(u) = f_{\beta}(b(u))\) .

The map \(f \mapsto f_{\beta}\) is a morphism of complex algebras from \(\mathcal{K}(\mathrm{Sp}(u))\) into \(\mathcal{C}(\mathrm{Sp}(b(u)))\) , so the map \(f \mapsto f(u)\) is a morphism of complex algebras from \(\mathcal{K}(\mathrm{Sp}(u))\) into \(\mathcal{L}(\mathrm{E})\) . We have \(\overline{f}(u) = f(u)^*\) since \(\overline{f}_{\beta} = \overline{f_{\beta}}\) . If \(f \geqslant 0\) then \(f_{\beta} \geqslant 0\) , hence \(f(u)\) is a positive endomorphism of E.

Remark. --- Suppose that u is a normal endomorphism of E, so that \(\mathrm{Sp}(u)\) is a compact subset of C. Since the bounded transform \(b(u)\) coincides with \(\beta(u)\) (formula (5)), and since \(\mathrm{Sp}(b(u)) = \beta(\mathrm{Sp}(u))\) (cor. 2 of prop. 7 of I, p. 111), the spectrum of \(b(u)\) is a compact subset of D. For every continuous function \(f \in \mathcal{C}(\mathrm{Sp}(u))\) , the function \(f_{\beta}\) coincides with the restriction of \(f \circ \beta^{-1}\) to \(\mathrm{Sp}(u)\) . Consequently, \(f_{\beta}(\beta(u))\) coincides with the endomorphism \(f(u)\) defined by the continuous functional calculus of u. The above definition is therefore compatible with that of the continuous functional calculus of the star algebra \(\mathcal{L}(\mathrm{E})\) .

Let \(f\in \mathcal{K}(\mathrm{Sp}(u))\) . We have

\[
\| f (u) \| \leqslant \| f _ {\beta} \| _ {\infty} = \| f \| _ {\infty},
\]

so that, for all \(x\) and \(y\) in E, we obtain the majorization \(|\langle x | f(u)y\rangle| \leqslant \|x\|\|y\|\|f\|_{\infty}\) . The map \(f \mapsto \langle x | f(u)y\rangle\) is therefore a bounded measure on \(\mathrm{Sp}(u)\) of total mass \(\leqslant \|x\|\|y\|\) (INT, IV, p. 154, § 4, no. 7). If \(x = y\) it is a positive measure, since \(f(u)\) is positive when \(f \geqslant 0\) .

DEFINITION 4. --- Let u be a normal partial operator on a complex Hilbert space E. Let x and y be in E. The bounded measure on Sp(u) defined by \(f \mapsto \langle x | f(u)y \rangle\) for \(f \in \mathcal{K}(\mathrm{Sp}(u))\) is called the spectral measure of \((x,y)\) relative to \(u\) . When \(x = y\) , we say that it is the spectral measure of \(x\) relative to \(u\) .

Remarks. --- 1) When u is continuous, this definition coincides with that of def. 2 of IV, p. 190.

\begin{enumerate}
\def\labelenumi{\arabic{enumi})}
\setcounter{enumi}{1}
\tightlist
\item
  Let \(j\) be the canonical inclusion of \(\mathrm{Sp}(u)\) into \(\mathbf{C}\) . Since \(\mathrm{Sp}(u)\) is closed, the map \(j\) is \(\mu\) -proper for every bounded measure \(\mu\) on \(\mathrm{Sp}(u)\) (INT, V, p. 68, § 6, no. 1, prop. 1). One often identifies the spectral measures of \(u\) with the measures on \(\mathbf{C}\) which are their images under \(j\) (cf. INT, V, p. 84, § 7, no. 2).
\end{enumerate}

The map from E×E into the Banach space \(\mathcal{M}^{1}(\mathrm{Sp}(u))\) that associates to \((x,y)\) the spectral measure of \((x,y)\) relative to u is sesquilinear and continuous with norm \(\leqslant 1\) .

Lemma 4. --- Let X be a locally compact topological space, let \(\mu\) be a positive measure on X and let \(g\) be a function \(\mu\)-measurable on X. Let \(m_g\) be the multiplication operator by \(g\) on \(\mathrm{L}^2 (\mathrm{X},\mu)\) .

\begin{enumerate}
\def\labelenumi{\alph{enumi})}
\item
  The application \(f\mapsto f(m_{g})\) from \(\mathcal{K}(\mathrm{Sp}(m_g))\) into \(\mathcal{L}(\mathrm{L}^{2}(\mathrm{X},\mu))\) is given by \(f\mapsto m_{f\circ g}\) ;
\item
  For \(f_1\) and \(f_2\) in \(\mathcal{L}^2(\mathrm{X},\mu)\) with classes \(\widetilde{f}_1\) and \(\widetilde{f}_2\) in \(\mathrm{L}^2(\mathrm{X},\mu)\) the spectral measure of \((\widetilde{f}_1,\widetilde{f}_2)\) relative to \(m_g\) is the restriction to \(\mathrm{Sp}(m_g)\) of the image measure \(g(\overline{f}_1f_2\cdot \mu)\) .
\end{enumerate}

The partial operator \(m_g\) is normal and \(b(m_g) = m_{\beta \circ g} = m_{g(1+|g|^2)^{-1/2}}\) (lemma 2). Moreover, \(\beta(g(x))\) belongs to \(\mathrm{Sp}(m_{\beta \circ g})\) for all \(x\) outside a locally \(\mu\) -negligible \(Y \subset X\) (proposition 22 of IV, p. 252 and lemma 1 of IV, p. 181). Since \(f_\beta(\beta(g(x))) = f(g(x))\) for all \(x \in X - Y\) , we deduce that \(f(m_g) = m_{f \circ g}\) by definition 3 and the corollary of Proposition 6 of IV, p. 187.

Let \(f_1\) and \(f_2\) be in \(\mathcal{L}^2(\mathrm{X},\mu)\) with classes \(\widetilde{f}_1\) and \(\widetilde{f}_2\) in \(\mathrm{L}^2(\mathrm{X},\mu)\) . Since the measure \(\nu = \overline{f}_1f_2\cdot \mu\) is bounded, the image measure \(g(\nu)\) is defined (INT, V, p. 69, § 6, no. 1, remark 1). Let \(f\in \mathcal{K}(\mathrm{Sp}(m_g))\) . We have \(f(m_{g}) = m_{f\circ g}\) according to \(a\) ), whence

\[
\langle \widetilde {f} _ {1} \mid f (m _ {g}) \widetilde {f} _ {2} \rangle = \int_ {\mathrm{X}} \overline {{f}} _ {1} (f \circ g) f _ {2} d \mu = \int_ {\mathrm{X}} (f \circ g) (\overline {{f}} _ {1} f _ {2} d \mu) = \int_ {\mathbf {C}} f d \nu
\]

(INT, V, p. 69, § 6, n° 1, formula (1)), which establishes assertion b).

Example. --- Let \(n \in \mathbf{N}\) . We equip \(\mathbf{R}^{n}\) with Lebesgue measure, and we identify \(\mathbf{R}^{n}\) and its dual group as in Corollary 3 of II, p. 236. Let N denote the Euclidean norm on \(\mathbf{R}^{n}\) and let \(\mathcal{F}\) be the Fourier transform on \(\mathcal{S}(\mathbf{R}^{n})\) and on \(\mathcal{S}'(\mathbf{R}^{n})\) (no. 12 of IV, p. 217).

Let \(\Delta_{\mathcal{S}}\) be the partial operator on \(\mathrm{L}^2 (\mathbf{R}^n)\) whose domain is the space \(\mathcal{S}(\mathbf{R}^n)\) and which satisfies

\[
\Delta_ {\mathcal {S}} (\varphi) = - \sum_ {i = 1} ^ {n} \partial_ {i} ^ {2} \varphi
\]

for all \(\varphi\in\mathcal{S}(\mathbf{R}^{n})\) . We then have \(\mathcal{F}(\Delta_{\mathcal{S}}(\varphi))=4\pi^{2}\mathrm{N}^{2}\mathcal{F}(\varphi)\) for every function \(\varphi\in\mathcal{S}(\mathbf{R}^{n})\) (remark 12 of IV, p. 220). Since the Fourier transform is an automorphism of \(\mathcal{S}(\mathbf{R}^{n})\) (prop. 18 of IV, p. 219), this means that the partial operator \(u=\mathcal{F}\circ\Delta_{\mathcal{S}}\circ\mathcal{F}^{-1}\) is the reduction of the multiplication operator \(m_{4\pi^{2}\mathrm{N}^{2}}\) on the space \(\mathcal{S}(\mathbf{R}^{n})\) . The partial operator \(u\) is closable and symmetric; its closure is the positive self-adjoint partial operator \(m_{4\pi^{2}\mathrm{N}^{2}}\) (prop. 6 of IV, p. 232, applied to the space \(\mathcal{D}(\mathbf{R}^{n})\subset\mathcal{S}(\mathbf{R}^{n})\) using prop. 4 of IV, p. 202). Consequently, the partial operator \(\Delta_{\mathcal{S}}\) is essentially self-adjoint. We denote \(\Delta\) its closure; it is a positive self-adjoint operator, and it is the unique Laplacian on \(\mathbf{R}^{n}\) ( \(cf.\) example from No. 6 of IV, p. 243 and corollary of Prop. 26 of IV, p. 257).

By prop. 21 of IV, p. 223, the Sobolev space \(\mathrm{H}^2 (\mathbf{R}^n)\) is the set of \(f\in \mathrm{L}^2 (\mathbf{R}^n)\) such that \((1 + \mathrm{N}^2)\mathcal{F}(f)\) belongs to \(\mathrm{L}^2 (\mathbf{R}^n)\), that is, such that \(\mathcal{F}(f)\) belongs to the domain of \(m_{4\pi^2\mathrm{N}^2}\) . Since the Fourier transform is an isometric isomorphism of the space \(\mathrm{L}^2 (\mathbf{R}^n)\) onto itself (th. 1 of II, p. 215), the domain of the Laplacian \(\Delta\) is \(\mathrm{H}^2 (\mathbf{R}^n)\) . We have \(\Delta = \mathcal{F}^{-1}\circ m_{4\pi^2\mathrm{N}^2}\circ \mathcal{F}\) . In particular, the spectrum of \(\Delta\) is equal to \(\mathbf{R}_{+}\) .
% semantic-fragments: legacy-input-seam
\relax{}
